\chapter{CONCLUSÃO}
\label{chap:conclusao}

Este trabalho comparou a estrutura de cooperação que emerge quando classificadores são treinados com amostras reais com as estruturas que emergem quando os mesmos classificadores são treinados com amostras sintéticas. Essas amostras foram produzidas por modelos probabilísticos parametrizados a partir dos dados reais de treinamento. O problema foi estudado em três condições---\RealToReal{}, \GaussianToReal{} e \GMMToReal{}---avaliadas sobre o mesmo conjunto de teste real fixo, com enumeração de todos os subconjuntos não vazios de canais e crescimento incremental do tamanho amostral por classe.

A principal conclusão é que a preservação estrutural é multidimensional. Correlação de ranking, \WinnersAgreement{} e similaridade progressiva de \Nstar{} respondem a perguntas diferentes e não podem ser substituídas por uma única medida. Os experimentos apresentaram ranking e relações de vencedores frequentemente elevados, enquanto a preservação de \Nstar{} foi mais variável e sensível à dinâmica dos cruzamentos.

A primeira pergunta de pesquisa recebeu resposta predominantemente positiva. Em grande parte das configurações, as estruturas obtidas nas condições sintéticas reproduziram fortemente a ordenação dos subconjuntos observada na condição real, com correlações de Spearman elevadas no maior tamanho amostral. A segunda pergunta apresentou resultado semelhante: a \WinnersAgreement{} mostrou que a maioria das decisões par a par foi mantida, mesmo quando o primeiro subconjunto do ranking não coincidiu exatamente. Esse resultado evidencia que a identidade isolada da primeira posição é uma descrição insuficiente da eficiência multicanal.

A terceira pergunta, referente a \Nstar{}, produziu o quadro mais heterogêneo. Houve casos de preservação baixa, moderada e alta. O maior valor final auditado ocorreu no Air Quality com Gaussian Naive Bayes e GMM, enquanto o Random Forest apresentou a maior similaridade dentro do SUPPORT2. A métrica revelou informação não capturada pelo ranking e pela \WinnersAgreement{}: condições quase idênticas na ordenação final podem divergir quanto à posição das últimas transições de superioridade.

A evolução com $n$ mostrou maior instabilidade nos menores regimes amostrais e tendência de estabilização das métricas de ordenação em valores maiores, sem monotonicidade universal. A similaridade de \Nstar{} pode oscilar porque cada novo prefixo introduz cruzamentos, altera últimos eventos e modifica o conjunto de células comparadas. Portanto, o tamanho amostral não é apenas uma variável de precisão; ele é parte da própria estrutura de cooperação.

Os resultados também evidenciaram a dependência do classificador. O Gaussian Naive Bayes apresentou preservação quase integral do ranking e dos vencedores no Air Quality. O Random Forest, embora tenha alcançado perdas próximas de 0,455 no SUPPORT2, preservou uma estrutura densa de últimos cruzamentos com similaridade de 0,6799 para a Gaussiana única e 0,6340 para o GMM. A presença de mais de 11 mil cruzamentos definidos em cada condição afasta a interpretação de que o resultado decorra de uma matriz trivial. Ao mesmo tempo, o caso demonstra que desempenho preditivo absoluto e fidelidade estrutural são propriedades distintas.

A comparação entre geradores não produziu um vencedor universal. O GMM reproduziu visualmente melhor certas geometrias, e o treinamento com suas amostras melhorou a similaridade de \Nstar{} em várias configurações. O treinamento com amostras da Gaussiana única, entretanto, reproduziu ranking e vencedores de forma igual ou superior em alguns casos, inclusive no Occupancy, apesar de essas amostras serem visualmente mais grosseiras. O princípio resultante é de parcimônia: a complexidade adicional do GMM deve ser adotada quando traz ganho estrutural relevante para o objetivo da análise.

Do ponto de vista metodológico, o CoInfoSim amplia a linha do SLACGS em quatro direções principais: ancora a simulação em datasets reais, substitui dimensões prefixadas por todos os subconjuntos de canais, avalia treinamento sintético sobre teste real e separa componentes de fidelidade estrutural. A infraestrutura produz rankings, relações de vencedores, matrizes de \Nstar{}, critérios de parada, metadados e relatórios regeneráveis.

A motivação operacional foi retomada por meio de uma agenda de custo-eficiência. Os resultados não demonstram economia nem recomendam a remoção de sensores ou exames. Eles fornecem os objetos necessários para uma futura integração com custos específicos de domínio. Ranking e Winners Agreement podem apoiar a identificação de alternativas eficientes; \Nstar{} pode informar quando canais adicionais passam a oferecer vantagem dentro da grade observada.

As limitações impedem generalizações amplas. Foram utilizados três datasets binários de baixa dimensão, o Random Forest foi executado apenas no SUPPORT2, as métricas derivam da perda 0--1 e a projeção para $n>512$ não foi validada. Também é necessário investigar a sensibilidade a particionamentos, prevalência, sementes, métricas alternativas e geradores capazes de representar dependência temporal, deriva da distribuição e ausência de dados.

Como trabalhos futuros, destacam-se: inclusão de jitter e outros geradores parametrizados; execução de classificadores não lineares em mais cenários; intervalos de confiança das métricas estruturais; validação fora da faixa amostral; modelagem explícita de custos; identificação de fronteiras de Pareto; e substituição da enumeração exaustiva por busca estruturada em problemas de maior dimensão.

Em síntese, os experimentos mostram que amostras sintéticas não precisam reproduzir integralmente a geometria dos dados reais nem as perdas absolutas para serem úteis ao estudo da cooperação. Em determinadas combinações de dataset, gerador e classificador, o treinamento sintético preserva de forma consistente a organização das alternativas e parte relevante de sua dinâmica amostral. O \CoInfoSim{} estabelece, assim, uma base experimental para investigar quando essa preservação é suficiente para apoiar análises e, futuramente, o planejamento de sistemas de aquisição de informação.
